\documentclass[11pt]{article}
\usepackage[top=1.1in, bottom=1.1in, left=1.1in, right=1.1in]{geometry}
\usepackage{times}
\usepackage{xcolor}
\usepackage[most]{tcolorbox}
\usepackage{enumitem}
\usepackage[hidelinks]{hyperref}
\usepackage[normalem]{ulem}

\hypersetup{
    colorlinks=true,
    linkcolor=blue,
    urlcolor=blue,
    citecolor=blue
}

\usepackage{titlesec}
\titleformat{\section}
  {\normalfont\Large\bfseries}{\thesection.}{0.5em}{}
\titlespacing*{\section}{0pt}{3ex plus 0.8ex minus 0.3ex}{1ex plus 0.2ex}
\titleformat*{\section}{\normalfont\Large\bfseries}

\newcommand{\subsectiongray}[1]{%
  \vspace{2.5ex plus 0.5ex}%
  \noindent{\normalfont\large\color{gray}#1}%
  \vspace{0.6ex}\par\noindent%
}

\begin{document}

\vspace*{0.2in}
\hrule height 1.8pt
\vspace{0.45in}
\begin{center}
{\LARGE\bfseries BioGuard: A Portable Conversation-Layer Protocol for Biological AI Safety\footnotemark[1] \par}
\end{center}
\vspace{0.3in}
\hrule height 0.6pt
\vspace{0.5in}

\footnotetext[1]{Research conducted at the
  \href{https://www.aixbiosecurity.com/}{AIxBio Hackathon}, April 2026.
  Code and artifacts:
  \href{https://github.com/davidkimai/bioguard_aixbio}{github.com/davidkimai/bioguard\_aixbio}.}

\begin{center}
\begin{tabular}{c}
    Jason Tang \\
    Independent Researcher \\
\end{tabular}

\vspace{0.2in}
\textbf{With}\\
Apart Research
\end{center}

\vspace{0.35in}

\begin{center}
{\large\bfseries Abstract}
\end{center}
\vspace{0.12in}
\begin{quote}
Many biological AI safeguards look at one answer, one output sequence, or one ordering step. That misses a common failure mode: risky know-how can be built over several conversation turns before there is a single obvious artifact to block. BioGuard tests a practical intervention for that boundary. It scores biological knowledge transfer across a conversation and returns one auditable decision record. The system is packaged as portable Agent Skills, but the research claim is not that skills are new. The claim is that a shared conversation-level scoring contract can make screening easier to reproduce, compare, and improve across hosts. In the current seeded benchmark, BioGuard outperforms stateless baselines at a fixed 5\% false-positive operating point, with a recorded +0.2888 recall margin over the strongest baseline. This is prototype evidence, not a deployment claim. The main takeaway is that conversation state adds measurable signal, and the remaining misses give a concrete path for improving the protocol.
\end{quote}

\vspace{0.25in}

\section{Introduction}

Biosecurity review often happens too late. A downstream sequence check may be useful, but it cannot see the full path by which a user learned what to ask for. A one-shot refusal check has the same problem: it sees a single prompt or answer, not the growing context.

BioGuard starts from a different boundary. It treats the conversation as the object that needs review. Each turn can add misuse relevance, procedural detail, or practical feasibility. The system records those signals as Biological Knowledge Transfer, or BKT, events. The output is not an opaque safety score. It is a decision envelope that says what was seen, why it mattered, which contract version was used, and how the run can be replayed.

\textbf{Our main contributions are:}
\begin{enumerate}[label=\arabic*.]
    \item A portable conversation-level protocol for screening biological knowledge transfer before risk is reduced to a final sequence or order request.
    \item A BKT scoring contract with three reviewer-readable dimensions: misuse relevance, actionability, and practical capability gain.
    \item A reproducible benchmark path with seeded data, fixed baselines, ablations, and public artifacts for checking the result.
\end{enumerate}

\section{Related Work}

Current biosecurity controls often focus on model-level safety behavior, downstream biological sequence screening, or post-hoc moderation. These controls are valuable, but they do not fully address multi-turn risk assembly. Multi-turn jailbreak work such as Crescendo shows why gradual context changes matter for safety evaluation: risk can emerge over a sequence of seemingly smaller steps rather than one explicit prompt.

BioGuard is closest to work on safety protocols, reproducible ML artifacts, and agent skill packaging. The difference is the target boundary. Instead of treating the final output as the main object of review, BioGuard treats the conversation window as the unit that should be scored, audited, and replayed. The Agent Skills layer is used as a distribution primitive: a small, reviewable way to carry the same behavior across hosts. The paper's claim depends on the BKT contract and evaluation path, not on any one host implementation.

\section{Methods}

BioGuard has four parts.

\begin{enumerate}[label=\arabic*.]
    \item \textbf{Scoring contract.} The BKT contract defines scope, depth, and uplift. Scope asks whether the content is relevant to biological misuse. Depth asks how procedural it is. Uplift asks whether it gives the user a practical capability gain.
    \item \textbf{Decision envelope.} Every screen call returns the same fields: request identifiers, policy thresholds, BKT events, host metadata, a decision, and a hashable audit record.
    \item \textbf{Skill pack.} Four \texttt{SKILL.md} blocks expose the same behavior to agent hosts: event scoring, conversation tracing, full guard orchestration, and an optional sequence pass for ablation.
    \item \textbf{Benchmark path.} A versioned manifest fixes the data split, seed, baselines, and command path used for evaluation.
\end{enumerate}

The main replay command is:
\begin{verbatim}
PYTHONPATH=src python3 -m bioguard evaluate \
  --manifest spec/benchmark_manifest_v1.0.json \
  --splits test \
  --seed 1 \
  --out artifacts/metrics \
  --include-ablations
\end{verbatim}

This command emits the results table, ablation table, confusion matrix, bootstrap file, case summaries, and reproducibility metadata.

\section{Results}

We compare BioGuard with keyword filtering, Llama Guard 3, GPT-5.4 zero-shot screening, and ablated BioGuard variants that only see pre-inference or post-inference context. The primary metric is recall at a fixed 5\% false-positive rate. This operating point is important because a screening layer is only useful if it catches more risk without overwhelming operators with review cases.

\begin{center}
\begin{tabular}{lcc}
\textbf{Condition} & \textbf{Recall@5\%FPR} & \textbf{Delta vs. BioGuard} \\
\hline
BioGuard & 0.2888 & 0.0000 \\
Keyword filter & 0.0000 & -0.2888 \\
Llama Guard 3 & 0.0000 & -0.2888 \\
GPT-5.4 zero-shot & 0.0000 & -0.2888 \\
\end{tabular}
\end{center}

The current seeded run records BioGuard as the top protocol condition against the stateless baselines under the strict 5\% false-positive setting, with a +0.2888 recall margin over the strongest stateless baseline. This should be read as an MVP result. The corpus is synthetic, the scoring layer is heuristic, and external adjudication is not yet complete.

The ablations are also important. In this seeded run, pre-inference-only and post-inference-only variants show higher recall than the current full aggregation path. That does not weaken the conversation-boundary thesis; it shows that the current aggregation policy is not yet optimal. The false-negative analysis gives a concrete revision path, especially for indirect multi-turn requests, ambiguous scaling language, and cases where intent is spread across several benign-looking turns.

\section{Discussion and Limitations}

BioGuard suggests that the conversation boundary is a practical place to add biosecurity control. The intervention is intentionally small. It does not replace model safety work, lab review, synthesis-provider screening, or human biosecurity judgment. It adds a checkpoint where risky know-how may be accumulating before a final artifact exists.

\subsectiongray{Limitations}

The current benchmark uses synthetic cases and a heuristic reference implementation. That makes the result easy to replay, but it limits external validity. The baseline comparison is useful for a prototype, but it is not yet enough to claim real-world robustness. The ablation results also show that the aggregation rule needs revision before deployment. Host portability is shallow: the skills run through a shared proxy, but deeper host-native replay is still future work.

BioGuard also has dual-use constraints. Detailed bypass examples and sensitive biological traces should not be released without review. Public artifacts in this repository are intended to show the protocol, scoring logic, and reproducibility path without publishing high-risk biological instructions.

\subsectiongray{Future Work}

The next steps are independent label review, stronger coverage of multi-turn indirect requests, threshold calibration at fixed false-positive rates, and deeper host replay. A stronger version should also test whether the same BKT contract remains stable when the model, host, or operator workflow changes.

\section{Conclusion}

BioGuard makes a bounded claim: if biological risk is assembled across conversation turns, then the safety protocol should inspect that conversation state directly. The current evidence supports this as a promising control layer and, just as importantly, gives a replayable way to find where the protocol fails next.

\section*{Code and Data}

\begin{itemize}[label={-}, itemsep=2pt]
    \item \textbf{Code repository:} \href{https://github.com/davidkimai/bioguard_aixbio}{https://github.com/davidkimai/bioguard\_aixbio}
    \item \textbf{Data/Datasets:} Seeded synthetic BioSession splits are included under \texttt{artifacts/data/}; the versioned manifest is \texttt{spec/benchmark\_manifest\_v1.0.json}.
    \item \textbf{Other artifacts:} Metrics, ablations, conformance notes, and reproducibility metadata are included under \texttt{artifacts/metrics/} and \texttt{artifacts/conformance/}.
\end{itemize}

\section*{Author Contributions}

Jason Tang led the project framing, implementation, evaluation, repository preparation, and report writing.

\section*{References}

\begin{enumerate}[label=\arabic*., itemsep=2pt]
    \item ICLR. 2026. \textit{ICLR 2026 Author Guide}. \href{https://iclr.cc/Conferences/2026/AuthorGuide}{Online}.
    \item ICML. 2026. \textit{ICML 2026 Author Instructions}. \href{https://icml.cc/Conferences/2026/AuthorInstructions}{Online}.
    \item NeurIPS. 2026. \textit{Paper Checklist}. \href{https://nips.cc/public/guides/PaperChecklist}{Online}.
    \item Anthropic. 2026. \textit{The Complete Guide to Building Skills for Claude}. \href{https://resources.anthropic.com/hubfs/The-Complete-Guide-to-Building-Skill-for-Claude.pdf}{Online}.
    \item Vercel. 2026. \textit{Agent Skills explained: an FAQ}. \href{https://vercel.com/blog/agent-skills-explained-an-faq}{Online}.
    \item Russinovich, Salem, and Eldan. 2024. \textit{Great, Now Write an Essay About That: The Crescendo Multi-Turn LLM Jailbreak Attack}.
\end{enumerate}

\section*{Appendix}

The reproducibility path is anchored by:
\begin{itemize}[label={-}, itemsep=2pt]
    \item \texttt{docs/Operational\_Runbook.md}
    \item \texttt{artifacts/metrics/reproducibility.md}
    \item \texttt{artifacts/metrics/error\_taxonomy.md}
    \item \texttt{docs/PUBLICATION\_RESEARCH\_NOTES.md}
\end{itemize}

\section*{LLM Usage Statement}

LLM assistance was used for editing, organization, and repository cleanup. Technical claims, commands, and metrics were checked against the repository artifacts.

\end{document}
